\documentclass[11pt]{article}

% ---------- Paquetes ----------
\usepackage[spanish]{babel}
\usepackage[utf8]{inputenc}
\usepackage[T1]{fontenc}
\usepackage{lmodern}
\usepackage[margin=2.5cm]{geometry}
\usepackage{xcolor}
\usepackage{titlesec}
\usepackage{enumitem}
\usepackage{amsmath}
\usepackage{parskip}
\usepackage[hidelinks,colorlinks=true,linkcolor=black,urlcolor=black]{hyperref}

% ---------- Colores ----------
\definecolor{bonsaigreen}{HTML}{2E5339}
\definecolor{notegray}{HTML}{808080}
\definecolor{placeholdergray}{HTML}{BFBFBF}
\definecolor{subgray}{HTML}{444444}

% ---------- Estilo de títulos ----------
\titleformat{\section}{\color{bonsaigreen}\normalfont\Large\bfseries}{\thesection.}{0.5em}{}
\titleformat{\subsection}{\color{subgray}\normalfont\large\bfseries}{\thesubsection}{0.5em}{}
\titlespacing*{\section}{0pt}{18pt}{8pt}
\titlespacing*{\subsection}{0pt}{14pt}{6pt}

\setlength{\parskip}{6pt}
\setlength{\parindent}{0pt}

% ---------- Comandos auxiliares ----------
% Nota guía: texto itálico gris, se borra al escribir el contenido final
\newcommand{\nota}[1]{{\itshape\color{notegray}\small NOTA: #1}\par\vspace{2pt}}
% Lista de puntos dentro de una nota guía
\newenvironment{notalist}{\begin{itemize}[itemsep=1pt,topsep=2pt,leftmargin=1.4em]\itshape\color{notegray}\small}{\end{itemize}\vspace{2pt}}
% Placeholder de contenido a completar
\newcommand{\placeholder}{\color{placeholdergray}[Escribir contenido acá]\color{black}\par\vspace{6pt}}

\begin{document}

% =========================================================
% PORTADA
% =========================================================
\begin{titlepage}
\centering
\vspace*{4cm}
{\color{bonsaigreen}\bfseries\large BONSAI CORP}\\[0.6cm]
{\bfseries\Huge Optimización de Packaging y Logística}\\[0.4cm]
{\color{notegray}\Large Informe Final --- Competencia de Datos AlixPartners}\\[2cm]

\large
Integrantes: Demián Elnecavé, Denis Wu, Emiliana Verdun\\[0.3cm]
Fecha de entrega: [dd/mm/aaaa]\\[0.3cm]
\nota{\color{red}ACORDARSE DE COMPLETAR}
Equipo en Kaggle: Zero Degree\\[0.3cm]

\end{titlepage}

% =========================================================
% ÍNDICE
% =========================================================
\tableofcontents
\nota{El índice se genera y actualiza automáticamente al compilar (puede requerir compilar dos veces para que los números de página queden correctos).}
\newpage

% =========================================================
% 1. RESUMEN EJECUTIVO
% =========================================================
\section{Resumen ejecutivo}

\nota{Media página máximo. Es lo primero (y a veces lo único) que lee el evaluador --- tiene que poder entender el resultado sin leer el resto.}
\begin{notalist}
\item Ahorro total logrado (USD y \%) vs. situación actual, desglosado en packaging y flete.
\item Reducción de \texttt{N\_tipos} de caja (de X a Y tipos).
\item 2--3 bullets con las decisiones clave que explican el resultado.
\item Una frase que mencione que se respetaron todas las restricciones (headspace, resistencia, fit, un producto por pallet).
\end{notalist}
\placeholder

% =========================================================
% 2. CONTEXTO Y OBJETIVO DEL DESAFÍO
% =========================================================
\section{Contexto y objetivo del desafío}

\subsection{El problema de negocio}

Bonsai Corp distribuye brócoli congelado desde 5 plantas hacia 5 regiones, y comercializa
un portafolio amplio de alrededor de 500 variantes de producto. Históricamente, cada nuevo
lanzamiento incorporó su propia caja secundaria --el envase de cartón que agrupa las
unidades del producto y que luego se apila en pallets para su transporte--- sin evaluar si
alguna caja ya existente podía reutilizarse. Como resultado, el catálogo actual está
compuesto por cientos de tipos de caja distintos, muchos de ellos con dimensiones
prácticamente idénticas entre sí.

Esta fragmentación tiene un costo real y creciente: cada tipo de caja implica su propio
molde de fabricación, sus propias corridas de producción con el proveedor y su propio
stock en depósito. Además, al pedirse en volúmenes bajos, muchos tipos de caja no acceden
a los descuentos por volumen que ofrece el proveedor, lo que encarece el costo unitario de
packaging. El problema no es exclusivamente de packaging: la eficiencia con la que las
cajas se apilan en el pallet determina también cuántos pallets se necesitan para
transportar la producción anual, y por lo tanto el costo de flete. Sin un criterio de
consolidación, el catálogo seguirá creciendo con cada nuevo lanzamiento, y con él, los
costos asociados.

\subsection{Objetivo de la optimización}

El objetivo del proyecto es reducir la cantidad de tipos de caja del catálogo,
consolidando productos con dimensiones compatibles en un mismo tipo de envase secundario,
sin comprometer la integridad del producto ni la eficiencia del transporte. Formalmente,
se busca minimizar el costo total:

\begin{equation*}
C_{\text{total}} = C_{\text{packaging}} + C_{\text{flete}}
\end{equation*}

donde $C_{\text{packaging}}$ surge del precio por caja ---definido por el grosor de cartón
elegido y el volumen anual consolidado de cada tipo--- y $C_{\text{flete}}$ surge de la
cantidad de pallets necesarios para transportar la producción anual, según el costo por
pallet de cada ruta planta-región (USD 150 dentro de la región de origen, USD 500 hacia el
resto).

Esta consolidación no está exenta de tensiones: una caja diseñada para servir a varios
productos a la vez puede no ser la más eficiente en el uso del espacio del pallet para
cada uno de ellos en particular, lo que podría aumentar el costo de flete incluso mientras
se reduce el costo de packaging. Identificar y cuantificar ese trade-off entre packaging y
flete es, por lo tanto, parte central del objetivo del proyecto y se desarrolla en detalle
en la Sección~6. Adicionalmente, la solución debe ser extensible: tiene que definir un
protocolo claro para que, ante cada nuevo lanzamiento, se pueda decidir si conviene
asignar una caja ya existente o crear una nueva, evitando que el catálogo vuelva a
fragmentarse con el tiempo.

% =========================================================
% 3. ANÁLISIS EXPLORATORIO DE DATOS
% =========================================================
\section{Análisis exploratorio de datos}
\nota{Esta sección está pedida explícitamente en la consigna y pesa en el criterio de ``rigor del análisis exploratorio''. El objetivo es mostrar el punto de partida antes de optimizar.}

\subsection{Portafolio actual de cajas}
\nota{Cantidad de tipos de caja, cuántos SKUs por tipo, dispersión de dimensiones (H, W, L). Gráfico: histograma o scatter de dimensiones mostrando clusters de productos ``casi iguales'' con cajas distintas.}
\placeholder

\subsection{Distribución de volúmenes y costos}
\nota{Volumen anual por producto y por tipo de caja actual. ¿Cuántos productos caen en Tier 1 ($<$ 20.000 u.) pagando el markup del 10\%? Costo actual desglosado: packaging vs. flete, por planta/región.}
\placeholder

\subsection{Utilización actual del pallet (\texorpdfstring{$U_{\text{pallet}}$}{U\_pallet})}
\nota{Baseline de \texttt{U\_pallet} por región/planta con el layout actual (column stacking, 1200~$\times$~800~mm, 1800~mm de altura útil).}
\placeholder

\subsection{Utilización volumétrica de la caja (\texorpdfstring{$U_{\text{caja}}$}{U\_caja}) y headspace actual}
\nota{Gap entre envase primario y caja. Sirve para justificar dónde hay margen para reajustar dimensiones (hasta 10\%) sin violar el headspace máximo.}
\placeholder

\subsection{Principales hallazgos del diagnóstico}
\nota{Síntesis en 3--5 bullets de las oportunidades detectadas, que después se retoman en la metodología.}
\placeholder

% =========================================================
% 4. METODOLOGÍA
% =========================================================
\section{Metodología}

\subsection{Criterio de consolidación de productos}
\nota{Cómo se agruparon productos en menos tipos de caja (ej. clustering por dimensiones/volumen). Justificar por qué ese criterio y no otro, en términos de negocio.}
\placeholder

\subsection{Selección de grosor de cartón}
\nota{Cómo se eligió entre 3,0 / 4,5 / 5,0 mm por grupo consolidado, balanceando resistencia (ECT) vs. costo. Recordar: un único grosor uniforme para toda la solución, sin mezclas.}
\placeholder

\subsection{Validación de restricciones}
\nota{Explicar cómo se verificó cada restricción --- es crítico porque incumplir cualquiera da puntaje 0 sin parcialidad:}
\begin{notalist}
\item Headspace máximo (6\% / 8\% / 10\% según grosor, tope 40~mm).
\item Reajuste de dimensiones internas hasta 10\% (volumen $\geq$ producto).
\item Resistencia a compresión ($\text{carga\_max} \geq$ peso de capas superiores).
\item Un solo producto por pallet.
\item Caja no rotable (H fija como eje vertical).
\end{notalist}
\placeholder

\subsection{Cálculo de utilización de pallet y pricing por tier}
\nota{Cómo se calculó \texttt{cantidad\_cajas\_total} por pallet (column stacking) y cómo se aplicó el factor de volumen (tiers 1--5) al precio base según grosor.}
\placeholder

% =========================================================
% 5. RESULTADOS
% =========================================================
\section{Resultados}

\subsection{Reducción en tipos de caja}
\nota{\texttt{N\_tipos} actual vs. propuesto. Gráfico de barras simple (antes/después).}
\placeholder

\subsection{Costo total: packaging + flete}
\nota{Costo actual vs. propuesto, desglosado por componente. Gráfico tipo waterfall: costo actual $\rightarrow$ ahorro packaging $\rightarrow$ ahorro flete $\rightarrow$ costo final.}
\placeholder

\subsection{Utilización de pallet por región}
\nota{\texttt{U\_pallet} promedio por región, escenario actual vs. propuesto. Idealmente un gráfico de barras agrupadas por región.}
\placeholder

\subsection{Tabla resumen de asignación}
\nota{Tabla agregada (no las 500 filas) a modo de muestra representativa. La tabla completa (\texttt{codigo\_producto}, \texttt{caja\_grosor\_mm}, \texttt{caja\_exterior\_largo}, \texttt{caja\_exterior\_ancho}, \texttt{caja\_exterior\_alto}) va en el Anexo A o como archivo .csv adjunto/linkeado.}
\placeholder

% =========================================================
% 6. TRADE-OFF PACKAGING VS. FLETE
% =========================================================
\section{Trade-off packaging vs. flete}
\nota{Pedido explícitamente en la consigna y es un criterio de evaluación aparte (``coherencia y justificación de negocio... incluyendo el trade-off'').}
\begin{notalist}
\item Mostrar casos donde consolidar cajas mejoró packaging pero empeoró \texttt{U\_pallet} (o viceversa).
\item Explicar cómo se ponderó esa tensión para llegar a la decisión final.
\item Un gráfico (scatter costo packaging vs. costo flete, o waterfall) que ilustre el trade-off.
\end{notalist}
\placeholder

% =========================================================
% 7. PROTOCOLO PARA NUEVOS PRODUCTOS
% =========================================================
\section{Protocolo de asignación para nuevos productos}
\nota{Parte del ``output esperado'' de la consigna. No necesita más matemática, sí buen diseño de proceso --- acá se puede lucir la creatividad.}
\begin{notalist}
\item Diagrama de flujo o checklist: cómo evaluar si un producto nuevo encaja en una caja existente.
\item Criterios de compatibilidad dimensional, headspace y resistencia a verificar.
\item Cuándo se justifica crear un tipo de caja nuevo en vez de reutilizar una existente.
\end{notalist}
\placeholder

% =========================================================
% 8. LIMITACIONES Y PRÓXIMOS PASOS
% =========================================================
\section{Limitaciones y próximos pasos}
\nota{Corto y honesto. Ej.: supuesto de column stacking, un solo layout de pallet para las 5 plantas, valores discretos de ECT. Qué se podría refinar con más datos o tiempo.}
\placeholder

% =========================================================
% 9. CONCLUSIONES
% =========================================================
\section{Conclusiones}
\nota{3--5 líneas cerrando el impacto de negocio y el mensaje principal para la audiencia no técnica.}
\placeholder

\newpage

% =========================================================
% ANEXOS
% =========================================================
\section*{Anexos}
\addcontentsline{toc}{section}{Anexos}

\subsection*{Anexo A: Tabla completa de asignación}
\addcontentsline{toc}{subsection}{Anexo A: Tabla completa de asignación}
\nota{Tabla completa, una fila por \texttt{codigo\_producto}, con columnas exactas: \texttt{codigo\_producto}, \texttt{caja\_grosor\_mm}, \texttt{caja\_exterior\_largo}, \texttt{caja\_exterior\_ancho}, \texttt{caja\_exterior\_alto}. Puede ir como link al .csv si el archivo es muy extenso para el documento.}
\placeholder

\subsection*{Anexo B: Metodología técnica / código (opcional)}
\addcontentsline{toc}{subsection}{Anexo B: Metodología técnica / código (opcional)}
\nota{La entrega del código es opcional y no es criterio de evaluación en sí mismo, pero si se incluye debe estar documentado y ser reproducible.}
\placeholder

\subsection*{Anexo C: Supuestos y fuentes de datos}
\addcontentsline{toc}{subsection}{Anexo C: Supuestos y fuentes de datos}
\nota{Listado breve de supuestos tomados sobre los 4 archivos de datos (catálogo de productos, especificaciones de cajas, operaciones de planta, compras).}
\placeholder

\end{document}
